\rowcolors{0}{}{}
\arrayrulecolor{black}
\renewcommand{\arraystretch}{1.3}
\small

\begin{longtable}{|C{1cm}|m{6.5cm}|m{7cm}|C{1cm}|}
\hline
\textbf{ID US} & \textbf{User Story} & \textbf{Description technique} & \textbf{Effort} \\
\hline
\endfirsthead

    \hline

\endhead

\hline
\endfoot

\hline
\endlastfoot

% Epic 1: Authentification et Sécurité
\multicolumn{4}{|l|}{\cellcolor{blue!20}\textbf{Epic 1 - Authentification et Sécurité}} \\
\hline

US01 & En tant qu'\textbf{utilisateur}, je souhaite m'authentifier avec email et mot de passe afin d'accéder à mes informations de manière sécurisée. & 
\vspace{2mm}
\begin{itemize}[nosep, left=0pt, label={$\bullet$}]
\item Développer l'endpoint \texttt{POST /api/auth/login}
\item Mettre en place la gestion de session via \textbf{JWT}
    (access token 15 min, refresh token 7 jours stocké en base)
\item Réaliser l'interface connexion login
\item Consommer l'API d'authentification
\item Gérer les tentatives échouées (3 essais) avec verrouillage temporaire
\end{itemize} & 2 \\
\hline

US02 & En tant qu'\textbf{utilisateur}, je souhaite réinitialiser mon mot de passe afin de retrouver l'accès à mon compte. & 
\vspace{2mm}
\begin{itemize}[nosep, left=0pt, label={$\bullet$}]
\item Développer les endpoints
    \texttt{POST /api/auth/mot-de-passe-oublie} et
    \texttt{POST /api/auth/reinitialiser-mot-de-passe}
  \item Créer les formulaires « mot de passe oublié » et
    de réinitialisation côté Next.js
  \item Générer un token sécurisé (bcrypt) avec expiration
    de \textbf{1 heure}, stocké dans
    \texttt{TokenReinitialisationMotDePasse}
  \item Invalider les anciens tokens non utilisés à chaque
    nouvelle demande
  \item Configurer l'envoi d'un e-mail de réinitialisation
    et d'un e-mail de confirmation via \textbf{Nodemailer}
\end{itemize} & 2 \\
\hline

US03 & En tant qu'\textbf{utilisateur}, je souhaite m'authentifier via Google ou GitHub afin de me connecter rapidement sans créer de compte. & 
\vspace{2mm}
\begin{itemize}[nosep, left=0pt, label={$\bullet$}]
\item Intégrer les stratégies OAuth2 via
    \texttt{passport-google-oauth20} et \texttt{passport-github2}
    dans NestJS (\texttt{@nestjs/passport})
\item Implémenter les routes \texttt{GET /api/auth/google},\texttt{GET /api/auth/github} et leurs callbacks
  \item Si l'utilisateur existe en base et est actif :
    générer les tokens JWT et rediriger vers le frontend
    avec \texttt{?oauth=success}
\item Vérifier l'existence de l'utilisateur en base de données et une demande d'inscription OAuth est envoyée à l'admin
\item Implémenter les routes /auth/google, /auth/github et leurs callbacks, avec génération d'un JWT
\item Implémenter les boutons OAuth sur la page de login
\end{itemize} & 3 \\
\hline

US04 & En tant qu'\textbf{utilisateur}, je souhaite me déconnecter de la plateforme afin de sécuriser ma session. & 
\vspace{2mm}
\begin{itemize}[nosep, left=0pt, label={$\bullet$}]
\item Développer l'endpoint POST /auth/logout avec révocation du refresh token en base
\item Supprimer \texttt{accessToken}, \texttt{refreshToken}
    et \texttt{user} du \texttt{localStorage}
\item Vider le cache React Query (\texttt{queryClient.clear()})
\item Réaliser le bouton de déconnexion 
\end{itemize} & 1 \\
\hline

US05 & En tant qu'\textbf{utilisateur}, je souhaite que mon compte soit protégé contre les tentatives de connexion malveillantes afin de garantir la sécurité de mes données. & 
\vspace{2mm}
\begin{itemize}[nosep, left=0pt, label={$\bullet$}]
  \item Implémenter le compteur de tentatives échouées
    (\texttt{tentativesConnexionEchouees}) en base de données
  \item Verrouiller le compte après \textbf{3 tentatives échouées}
    pendant \textbf{1 minute} via le champ
    \texttt{compteVerrouilleJusquA}
  \item Réinitialiser le compteur à 0 après une connexion réussie
  \item Appliquer le rate limiting global via
    \texttt{@nestjs/throttler} sur les endpoints publics
\end{itemize} & 2 \\
\hline

% Epic 2: Gestion des utilisateurs
\multicolumn{4}{|l|}{\cellcolor{blue!20}\textbf{Epic 2 - Gestion des utilisateurs}} \\
\hline

US06 & En tant qu'\textbf{utilisateur}, je souhaite gérer mon profil afin que mes informations soient toujours à jour. & 
\vspace{2mm}
\begin{itemize}[nosep, left=0pt, label={$\bullet$}]
\item Développer les endpoints \texttt{GET /api/profil} et \texttt{PUT /api/profil}
\item Implémenter l'upload de la photo de profil \texttt{POST /api/profil/upload-photo}
\item Développer le changement de mot de passe via
    \texttt{POST /api/profil/change-password}
    avec vérification de l'ancien mot de passe (bcrypt)
\item Gérer les e-mails secondaires :
    \texttt{GET/POST/DELETE /api/profil/emails} et
    \texttt{PATCH /api/profil/emails/principal}
\item Réaliser l'interface profil côté Next.js  avec formulaire prérempli
\end{itemize} & 2 \\
\hline

US07 & En tant qu'\textbf{administrateur}, je souhaite créer des utilisateurs afin de gérer les accès au système. & 
\vspace{2mm}
\begin{itemize}[nosep, left=0pt, label={$\bullet$}]
\item Développer l'endpoint \texttt{POST /api/users}
\item Générer automatiquement un mot de passe temporaire haché via \textbf{bcrypt} (facteur 10)
\item Générer un token de réinitialisation (1 heure) et construire le lien de réinitialisation
\item Créer automatiquement le \texttt{ProfilEmploye}
    avec matricule \texttt{EMP-\{id\}} si le rôle n'est pas CLIENT
\item Valider les données d'entrée via \textbf{class-validator}
  \item Envoyer un e-mail de bienvenue via \textbf{Nodemailer}
    contenant le mot de passe temporaire et le lien de
    réinitialisation
  \item Réaliser le formulaire d'ajout côté Next.js
\end{itemize} & 3 \\
\hline

US08 & En tant qu'\textbf{administrateur}, je souhaite modifier les informations d'un utilisateur afin de maintenir des données à jour. & 
\vspace{2mm}
\begin{itemize}[nosep, left=0pt, label={$\bullet$}]
\item Développer l'endpoint \texttt{PUT /api/users/:id}
\item Sécuriser l'accès à l'endpoint pour les administrateurs uniquement
\item Valider les données via \textbf{class-validator}
\item Réaliser le formulaire de modification côté Next.js
\end{itemize} & 2 \\
\hline

US09 & En tant qu'\textbf{administrateur}, je souhaite activer et désactiver un compte utilisateur afin de contrôler l'accès au système. & 
\vspace{2mm}
\begin{itemize}[nosep, left=0pt, label={$\bullet$}]
\item Développer les endpoints \texttt{PATCH /api/users/:id/desactiver} et \texttt{PATCH /api/users/:id/reactiver}
\item Implémenter les vérifications de réassignation des tickets et des demandes en cas de désactivation
\item Envoyer des notifications pour informer l'utilisateur de l'activation ou désactivation de son compte
\item Réaliser les boutons d'action avec appels des API
\end{itemize} & 3 \\
\hline

US10 & En tant qu'\textbf{administrateur}, je souhaite consulter la liste des utilisateurs afin de visualiser et gérer efficacement les comptes. & 
\vspace{2mm}
\begin{itemize}[nosep, left=0pt, label={$\bullet$}]
\item Développer l'endpoint \texttt{GET /api/users}
\item Réaliser l'interface d'administration avec vue grille
    et vue liste, filtres par statut et par rôle,
    barre de recherche dynamique

\end{itemize} & 1 \\
\hline

US11 & En tant qu'\textbf{utilisateur}, je souhaite consulter l'historique de mes activités afin de suivre mes actions dans le système. & 
\vspace{2mm}
\begin{itemize}[nosep, left=0pt, label={$\bullet$}]
\item Développer l'endpoint \texttt{GET /api/users/:id/activities} pour récupérer l'historique
\item Implémenter la logique de filtrage dans le backend en utilisant l'opérateur contains
\item Implémenter une barre de recherche dynamique
\item Réaliser l'interface timeline avec filtres par date et type d'action
\end{itemize} & 2 \\
\hline

% Epic 3: Gestion des rôles et permissions
\multicolumn{4}{|l|}{\cellcolor{blue!20}\textbf{Epic 3 - Gestion des rôles et permissions}} \\
\hline

US12 & En tant qu'\textbf{administrateur}, je souhaite gérer des rôles afin de structurer les accès au système. & 
\vspace{2mm}
\begin{itemize}[nosep, left=0pt, label={$\bullet$}]
\item Développer les endpoints CRUD pour la gestion complète des rôles
\item Implémenter le PermissionsGuard et le décorateur @Permissions pour sécuriser les routes du backend 
\item Réaliser l'interface de gestion sous forme de matrice
\end{itemize} & 3 \\
\hline

US13 & En tant qu'\textbf{administrateur}, je souhaite attribuer des permissions granulaires aux rôles afin de contrôler précisément les accès. & 
\vspace{2mm}
\begin{itemize}[nosep, left=0pt, label={$\bullet$}]
\item Développer l’endpoint PATCH /roles-permissions/roles/:id/sync pour synchroniser dynamiquement les associations entre un rôle et ses permissions. 
\item Implémenter la logique de mise à jour atomique                                                                                                                                                                                                                       
\item Développer l'interface interactive de type "Matrice de contrôle d'accès" 

\end{itemize} & 2 \\
\hline
US14 & En tant qu'\textbf{administrateur}, je souhaite consulter la matrice rôles-permissions afin de visualiser rapidement les droits de chaque rôle. & 
\vspace{2mm}
\begin{itemize}[nosep, left=0pt, label={$\bullet$}]
\item Développer l’endpoint GET /roles-permissions/matrix ) pour récupérer les données 
\item Implémenter la logique de structuration des données au format JSON                
\item Réaliser l’interface de visualisation de type "Heatmap" ou "Tableau de bord"

\end{itemize} & 3 \\
\hline

% Epic 4: Demandes d'inscription OAuth
\multicolumn{4}{|l|}{\cellcolor{blue!20}\textbf{Epic 4 - Demandes d'inscription OAuth}} \\
\hline

US15 & En tant qu'\textbf{utilisateur externe}, je souhaite demander l'accès via OAuth afin de m'inscrire facilement. & 
\vspace{2mm}
\begin{itemize}[nosep, left=0pt, label={$\bullet$}]
\item Développer l'endpoint POST /auth/oauth/request pour créer une demande
\item Enregistrer la demande en base avec statut EN\_ATTENTE
\item Implémenter l'envoi de notification à l'administrateur via email et in-app
\item Réaliser l'interface de demande d'accès
\end{itemize} & 3 \\
\hline

US16 & En tant qu'\textbf{administrateur}, je souhaite valider ou rejeter les demandes d'inscription OAuth afin de contrôler les accès. & 
\vspace{2mm}
\begin{itemize}[nosep, left=0pt, label={$\bullet$}]
\item Développer les endpoints PATCH /auth/oauth/requests/:id/approve et /reject
\item Implémenter la création automatique du compte utilisateur lors de l'approbation
\item Configurer l'envoi d'emails de confirmation ou rejet via Nodemailer
\item Réaliser l'interface de gestion des demandes avec actions d'approbation/rejet
\end{itemize} & 3 \\
\hline
% Epic 4: Demandes d'inscription OAuth


\end{longtable}